\chapter{Заключение: единая картина}
\label{ch:conclusion}

\section{Один субстрат, много операторов}

Если из монографии стоит унести одну мысль, то вот она: память Рин — это
\emph{один типизированный темпоральный граф с метаданными}, а всё богатство
поведения — \emph{операторы} над ним. Retrieval, пропагация уверенности,
забывание и консолидация, ночные обходы, настроение и слой личности — не отдельные
подсистемы со своими моделями памяти, а разные способы читать и писать одни и те
же поля (\texttt{confidence}, \texttt{activation}, \texttt{valence},
\texttt{provenance}, \texttt{need\_to\_discover}, временной слой) на одном графе.

\section{Четыре кита связаны}

Отсюда — практический вывод о разработке. Четыре кита (правильный retrieval,
правильная пропагация уверенности, правильное забывание/консолидация, правильные
ночные обходы) сцеплены через общее состояние: забывание меняет топологию,
топология — retrieval, retrieval — что пропагируется, пропагация — что найдёт скраб.
Их нельзя тюнить порознь. Именно поэтому нужен стенд с ground truth
(глава~\ref{ch:experiment}): без него взаимодействие операторов не измерить, и
отладка «на живой Рин» найдёт проблемы слишком поздно.

\section{Сквозные принципы}

Несколько идей проходят через всю монографию и заслуживают повторения:

\begin{itemize}
\item \textbf{Гипотеза, а не факт.} Всё, что порождено выводом (fusion, реакции),
пишется провизорно, с низкой уверенностью и провенансом; управление — на границе
коммита. Это защищает от превращения обаятельной ассоциативности в фабрику ложных
убеждений.
\item \textbf{Внутрь, а не наружу.} Значимость (гейт консолидации характера)
берётся из само-оценивания, а не из внешнего одобрения; иначе строится угодник, что
противоречит ядру личности.
\item \textbf{Ингредиенты, а не реакции.} Граф хранит материал; реакция
композируется под настроение и мотивацию, а личность — смещение во время обхода, а
не таблица весов.
\item \textbf{Континуум, а не бинарность.} Врождённое/приобретённое,
семантическое/эпизодическое, стабильное/пластичное — это оси глубины консолидации,
а не выключатели.
\item \textbf{Топология незрима, но правит.} Любой локальный процесс на графе
наследует спектрально-топологический транспорт; телепорт-скраб — не украшение, а
необходимая эпистемическая гигиена с доказуемой гарантией покрытия.
\end{itemize}

\section{Единая спина двух проектов}

И последнее — то, ради чего стоило смотреть память рядом с DTI. Консолидация есть
огрубление графа; тот же вопрос «что выживает при переходе fine\,$\to$\,coarse»
стоит и там, и там. Over-squashing, спектр лапласиана, эффективное сопротивление,
персистентный лапласиан, обучаемые операторы передачи сообщений — инструменты,
работающие в обе стороны. Память Рин — «грязный», богатый испытательный стенд для
тех же операторов огрубления, где нет чистых меток; DTI — их «чистая» лаборатория с
измеримой истиной. Отрефлексировать это явно значит получить не два проекта, а один
с двумя приложениями.

\section{Что делать дальше}

Порядок разработки, вытекающий из монографии: сначала прибить \emph{схему
субстрата} (узел/ребро, метаданные, что значит коммит и кто его гейтит); затем
собрать Ярус-1 стенд и погнать батарею тестов, начиная с самого дешёвого
(концентрация $\pi$ на хабах — сразу показывает, что диагностика жива); затем
локальное правило $\Delta\text{conf}$ (знак, гейт источника, дисконт
независимости), поверх которого пороги и ветвящийся коэффициент навешиваются
тривиально; параллельно — построить $\varphi$-скорер как измерительный прибор для
отбора голоса и фильтрации выходов. Настроение, сон и слой личности — усилители
поверх субстрата, которому сначала надо стать доверенным.
